\section{Summing the costs and comparing dimensions}\label{sec:tree}
We now assemble the recursion. The purpose of the sparse estimate is to prevent costs from accumulating indefinitely on diagrams with a long row or column. The remaining costs can be summed before we compare them with the dimension of the root.

Write $\log^+x=\max\{0,\log x\}$. Fix a root $\lambda\vdash N$, with $K=k(\lambda)>0$. A node with diagram $\alpha\vdash m$ is expanded only if $m$ exceeds a fixed cutoff $M$ and $k(\alpha)>m^{.01}$. Otherwise it is stopped. Choose the even moment $r$ large enough for Lemma~\ref{lem:sparse} and, by Lemma~\ref{lem:gap}, large enough that $F(\alpha)\le1$ for every stopped diagram of size at most $M$. These choices are possible simultaneously. The numerical cutoff will be fixed below before $r$ is finally chosen.

Squaring \eqref{eq:recursion} and bounding a sum by its number of terms times its largest term gives
\begin{equation}\label{eq:log-recursion}
 \log F(\alpha)\le2\log|\mathcal I(\alpha)|
 +\max_{\boldsymbol\beta\in\mathcal I(\alpha)}
     \sum_{i=1}^4\{\log F(\beta_i)+\theta\log C(\beta_i)\}.
\end{equation}
We use $\log0=-\infty$. If $F(\alpha)=0$, no bound is needed. Iterating this inequality amounts to taking the maximum over finite four-ary trees of occurring restrictions. Each expanded node contributes its choice cost $2\log|\mathcal I(\alpha)|$, and each of its four immediate children contributes $\theta\log C(\beta_i)$. The latter charge remains present even if the child is stopped.

The tail count in the proof of Lemma~\ref{lem:restriction-cost} gives
\[
 \log|\mathcal I(\alpha)|
 \le C(1+\sqrt{k(\alpha)})\log(2m).
\]
To apply it, orient the parent along its longest side, use the first-row inequality, and count each child's tail with at most the parent's deficit. Child sizes are prescribed. Thus this choice cost fits within the lower-order term of \eqref{eq:restriction-cost}. We may charge it at the parent, including the root, without altering the estimates below.

For $H>2$, put
\[
 R_H(\lambda)=\sum_{(i,j)\in\lambda}\log^+\frac{\min(i,j)}H.
\]
This quantity will count, box by box, the levels at which a root box
can lie outside the chosen hook. The budget below adds an entropy
term $\tfrac12K\log(N/K)$. Lemma~\ref{lem:hook-comparison} will
compare $R_H(\lambda)+K\log(N/K)$ with $\log D_\lambda$, with
explicit errors. We therefore seek a coefficient below one in front
of the budget, leaving room to absorb its lower-order costs.

\begin{lemma}[Uniform tree budget]\label{lem:tree-budget}
Fix $\delta\in(0,1/2)$ and put $a=1/2-\delta$. For every $H>2$ and $\epsilon>0$, a sufficiently large fixed cutoff $M$ makes every such recursion tree satisfy
\begin{equation}\label{eq:tree-budget}
 \log F(\lambda)\le\epsilon K+
 \frac\theta a\left\{
    R_H(\lambda)+\frac K2\log\frac NK\right\}.
\end{equation}
The cutoff and all constants in this statement are independent of the final even moment $r$.
\end{lemma}
\begin{proof}
At a level with node size $m$, let $q_m$ be the number of charged nodes and let their deficits be $k_1,\ldots,k_{q_m}$. By Lemma~\ref{lem:hook-deficit}, their sum is at most $K$. Expanded nodes satisfy $k_i>m^{.01}$, so their number is at most $K/m^{.01}$. A charged child of size $m$ has an expanded parent of size $4m$, and there are four children per parent. Consequently, after changing one absolute constant,
\begin{equation}\label{eq:charged-count}
 q_m\le C K m^{-.01},\qquad q_m\le N/m,
 \qquad\sum_i k_i\le K.
\end{equation}
The root also satisfies these bounds when expanded. A stopped root already has $F\le1$ and satisfies the conclusion directly.

At that level choose
\[
 h_m=\left\lceil m^a\sqrt{K/N}\right\rceil.
\]
The lower-order terms in the choice and restriction costs are bounded by
\[
 C\sum_i(1+h_m)(1+\sqrt{k_i})\log(2m).
\]
Cauchy--Schwarz gives $\sum_i\sqrt{k_i}\le\sqrt{q_mK}$. The terms arising from the ceiling and the constant one are therefore at most
$CKm^{-.005}\log(2m)$. For the terms involving $m^a$, use both bounds in \eqref{eq:charged-count}:
\[
 m^a\sqrt{K/N}\sqrt{q_mK}\le Km^{-\delta},
\]
\[
 m^a\sqrt{K/N}\,q_m
 \le CKm^{-\delta-.005}.
\]
For the second inequality, $q_m\le\sqrt{(N/m)(CK/m^{.01})}$. These costs are summable geometric tails because level sizes differ by a factor four. All charged nodes have size greater than $M/4$, including the stopped children of the last expanded level. Choosing $M$ large makes their total at most $\epsilon K$.

It remains to sum the leading hook terms. At every level, iteration of Lemma~\ref{lem:hook-deficit} bounds the sum of $t_{h_m}$ over charged nodes by $t_{h_m}(\lambda)$. Follow one root box $(i,j)$ and write $s=\min(i,j)$. It can contribute only if
\[
 s>h_m\ge m^a\sqrt{K/N}.
\]
The eligible sizes $m$ form part of a geometric progression of ratio four, all larger than $M/4$. By increasing $M$ so that $(M/4)^a\ge4^a H$, their number is at most
\[
 \frac{\log^+(s/H)+\tfrac12\log(N/K)}{a\log4}.
\]
Indeed the number is bounded by the positive part of one plus the logarithm of the ratio of the upper endpoint
$(s\sqrt{N/K})^{1/a}$ to the lower endpoint $M/4$, divided by $\log4$; the choice of $M$ absorbs the extra one. Replacing $\log(s/H)$ by its positive part can only enlarge the bound.

Boxes in the longest first row or column never contribute, since $h_m\ge1$. There are exactly $K$ remaining boxes. Multiply their level counts by $\theta\log4$ and add the lower-order costs. This proves \eqref{eq:tree-budget}.
\end{proof}

We have reduced the operator recursion to a purely diagrammatic comparison. The next lemma compares $R_H(\lambda)$ with the hook-length formula, retaining the first-row entropy needed when $K$ is small relative to $N$.

\begin{lemma}[Hook comparison]\label{lem:hook-comparison}
Orient $\lambda\vdash N$ so its first row is a longest side, put $K=N-\lambda_1>0$, and let $\zeta=(\lambda_2,\lambda_3,\ldots)\vdash K$. For $H>2$,
\begin{align}
 \log D_\lambda
 &\ge K\log(N/K)+\log D_\zeta-E_{N,K},\label{eq:row-hook}\\
 R_H(\lambda)&\le\log D_\zeta+4K/H,\label{eq:rank-hook}
\end{align}
where $E_{N,K}=O(K/N)$ for $K<N/3$ and
$E_{N,K}=O(\sqrt N\log(2N))$ for $K\ge N/3$.
There is also an absolute $c>0$ such that
$\log D_\lambda\ge cK$ for all sufficiently large $N$ and every nontrivial, nonsign $\lambda$.
\end{lemma}
\begin{proof}
The hook-length formula separates the first row exactly:
\[
 D_\lambda=\binom NKD_\zeta
 \prod_{j=1}^{\lambda_1}
 \left(1+\frac{\lambda'_j-1}{\lambda_1-j+1}\right)^{-1}.
\]
If $K<N/3$, a nonzero numerator occurs only for $j\le K$, so the logarithm of the inverted product is at most $K/(N-2K+1)$. For $K\ge N/3$, split the columns at height $\sqrt N$. There are at most $\sqrt N$ columns taller than this; each contributes at most $\log(1+N)$. In the remaining columns, bound $\log(1+x)$ by $x$ and sum the denominators harmonically. Their contribution is at most $\sqrt N(1+\log N)$. Since $\binom NK\ge(N/K)^K$, this proves \eqref{eq:row-hook}.

Every box contributing to $R_H$ lies in $\zeta$. At root coordinates $(i,j)$, its hook in $\zeta$ has length at most
\[
 \frac K{i-1}+\frac Kj\le\frac{4K}{\min(i,j)}.
\]
Take a reverse linear extension of the Young diagram order on $\zeta$, numbering its boxes from $1$ to $K$. The number assigned to each box is at least its hook length, because all boxes to its right and below precede it. The product of all number-to-hook ratios is exactly $D_\zeta$, and every ratio is at least one.

Suppose $t$ boxes contribute to $R_H$. Their distinct assigned numbers have product at least $t!$. Keeping only their ratios in the hook product gives, for $t>0$,
\[
 \log D_\zeta\ge R_H-t\log\frac{4eK}{tH}
 \ge R_H-\frac{4K}H.
\]
The last inequality follows by maximizing $t\log(4eK/(tH))$ over $t>0$. For $t=0$, the assertion is immediate.

Finally we prove $\log D_\lambda\ge cK$ without losing uniformity in the shape. If $K<N/3$, \eqref{eq:row-hook} already gives this for large $N$. Put $L=\lambda_1=\max(\lambda_1,\lambda'_1)$. If $K\ge N/3$ and $L\ge N/8$, retain the first row and an order ideal of $b=\lfloor N/32\rfloor$ boxes in the lower diagram. Place $1,\ldots,b$ in the first $b$ boxes of the top row, choose any $b$ of the remaining $L$ labels for the lower diagram, fill those lower boxes in a fixed standard order, and fill the rest of the top row increasingly. This gives $\binom Lb\ge4^b$ standard tableaux of the retained diagram. Every such tableau extends to $\lambda$. If $L<N/8$, every hook is shorter than $N/4$, so
$D_\lambda\ge N!/(N/4)^N\ge(4/e)^N$.
These cases give one absolute positive $c$.
\end{proof}

\begin{proof}[Proof of the moment assertion in Theorem~\ref{thm:moment}]
Fix the constants in their logical order. Theorem~\ref{thm:permanent} first gives $p_0$ and $\theta<1/2$. Choose $\delta>0$ so small that
\[
 b:=\frac\theta{1/2-\delta}<1.
\]
The universal degree constant $c$ in Lemma~\ref{lem:hook-comparison} is independent of this choice. Choose $H$ sufficiently large and $\epsilon$ sufficiently small that
$\epsilon+4b/H<(1-b)c/4$. Then choose the cutoff $M$ required by Lemma~\ref{lem:tree-budget}, increasing it further so that the errors $E_{N,K}$ in Lemma~\ref{lem:hook-comparison} are less than $(1-b)cK/(4b)$ whenever $N>M$. This is possible uniformly: the first error divided by $K$ is $O(1/N)$, and the second is $O(\log N/\sqrt N)$ in its stated range.

Only now choose the even moment $r$. It must make the four-site column law admissible for Theorem~\ref{thm:permanent}, enforce the sparse stopping rule, and make $F_r\le1$ for all nontrivial shapes of all dyadic sizes at most $M$. These are finitely many or fixed-threshold requirements, by Lemmas~\ref{lem:gap} and \ref{lem:sparse}; all persist on increasing $r$.

For an expanded root, \eqref{eq:tree-budget}, \eqref{eq:row-hook} and \eqref{eq:rank-hook} give
\[
 \log F_r(\lambda)
 \le b\log D_\lambda+\epsilon K+4bK/H+bE_{N,K}
 \le\frac{1+b}{2}\log D_\lambda.
\]
Here we discarded the additional negative term
$-bK\log(N/K)/2$. Thus $\eta=(1-b)/2>0$ works at every expanded root. At a stopped root $F_r\le1\le D_\lambda^{1-\eta}$. The row representation has $F_r=1$, while the column representation is killed by Lemma~\ref{lem:gap}. This proves the moment assertion for every dyadic $N$.
\end{proof}
